\chapter{DL-002 --- Canonicalize Patterns and Blocks}\label{record-169}

\noindent{\footnotesize\color{BookMuted}\path{docs/design/DL-002-canonicalize-patterns-and-blocks.md}\par}\smallskip

Legend: \href{https://github.com/flyingrobots/bijou/blob/1853ad902f6bf219326409e7f743a5a1d6517dbc/docs/legends/DL-design-language.md}{DL --- Design Language}

\section{Why this cycle exists}\label{record-169-why-this-cycle-exists}

DL-001 captured the doctrine layer.

That was necessary, but it still leaves Bijou with a gap between:

\begin{itemize}
\tightlist
\item
  doctrine
\item
  component families
\item
  shell/product work
\end{itemize}

This cycle exists to bridge that gap by making the most proven Bijou interaction language explicit:

\begin{itemize}
\tightlist
\item
  selection treatment
\item
  active-region treatment
\item
  drawer and panel rhythm
\item
  preference-row/value layout
\item
  the difference between patterns and future blocks
\end{itemize}

\section{Scope of this cycle}\label{record-169-scope-of-this-cycle}

This cycle intentionally covers:

\begin{itemize}
\tightlist
\item
  documenting canonical patterns in the design-system layer
\item
  explicitly defining what Bijou means by a block
\item
  identifying the first credible block candidates
\item
  updating the Design System read order so patterns and blocks are both first-class
\item
  spawning the next backlog slice for proving these patterns in shared surfaces
\end{itemize}

It does \textbf{not} include:

\begin{itemize}
\tightlist
\item
  shipping a full blocks library
\item
  replacing the app frame with a new block system
\item
  inventing new shell behavior
\item
  reworking notification center polish, which remains parked on the Humane Terminal backlog
\end{itemize}

\section{Human users}\label{record-169-human-users}

\subsection{Primary human user}\label{record-169-primary-human-user}

A Bijou builder or maintainer trying to turn doctrine into reusable interface decisions.

They need to know:

\begin{itemize}
\tightlist
\item
  which interaction treatments are canonical
\item
  what visual language means selection versus focus versus support text
\item
  how future blocks should be composed without inventing a parallel design language
\end{itemize}

\subsection{Human hill}\label{record-169-human-hill}

A builder can point to concrete Bijou patterns and block rules instead of translating abstract doctrine into ad hoc component choices.

\section{Agent users}\label{record-169-agent-users}

\subsection{Primary agent user}\label{record-169-primary-agent-user}

An agent reviewing or generating Bijou UI work.

It needs to:

\begin{itemize}
\tightlist
\item
  identify whether a surface follows canonical selection and spacing rules
\item
  distinguish a reusable pattern from a future reusable block
\item
  know which next cycle should turn the documented patterns into more shared code
\end{itemize}

\subsection{Agent hill}\label{record-169-agent-hill}

An agent can evaluate or propose Bijou surfaces against explicit pattern language instead of inferring those rules from scattered shell code.

\section{Human playback}\label{record-169-human-playback}

\begin{enumerate}
\def\labelenumi{\arabic{enumi}.}
\tightlist
\item
  A maintainer opens the Design System docs.
\item
  They can find an explicit explanation of canonical selection treatment, active-region treatment, drawer rhythm, and value-overflow behavior.
\item
  They can tell what counts as a block and why blocks come after patterns and components.
\item
  They can see the next backlog item for turning those documented patterns into more shared product surfaces.
\end{enumerate}

\section{Agent playback}\label{record-169-agent-playback}

\begin{enumerate}
\def\labelenumi{\arabic{enumi}.}
\tightlist
\item
  An agent reads the DL legend and DL-002 cycle doc.
\item
  It can locate the design-system docs that define canonical patterns and block rules.
\item
  It can determine whether selection, focus, and spacing treatments are now documented in one place.
\item
  It can point to the next backlog item for implementation-oriented follow-through.
\end{enumerate}

\section{Linked invariants}\label{record-169-linked-invariants}

\begin{itemize}
\tightlist
\item
  \href{https://github.com/flyingrobots/bijou/blob/cbd70612140e7a450d764b3c7be1c6af878b99bc/docs/invariants/tests-are-the-spec.md}{Tests Are the Spec}
\item
  \href{https://github.com/flyingrobots/bijou/blob/cbd70612140e7a450d764b3c7be1c6af878b99bc/docs/invariants/focus-owns-input.md}{Focus Owns Input}
\item
  \href{https://github.com/flyingrobots/bijou/blob/cbd70612140e7a450d764b3c7be1c6af878b99bc/docs/invariants/visible-controls-are-a-promise.md}{Visible Controls Are a Promise}
\item
  \href{https://github.com/flyingrobots/bijou/blob/cbd70612140e7a450d764b3c7be1c6af878b99bc/docs/invariants/graceful-lowering-preserves-meaning.md}{Graceful Lowering Preserves Meaning}
\item
  \href{https://github.com/flyingrobots/bijou/blob/cbd70612140e7a450d764b3c7be1c6af878b99bc/docs/invariants/docs-are-the-demo.md}{Docs Are the Demo}
\end{itemize}

\section{Invariants for this cycle}\label{record-169-invariants-for-this-cycle}

\begin{itemize}
\tightlist
\item
  background fill should remain the canonical selected-item treatment
\item
  structural accents such as gutters should remain the canonical active-region treatment
\item
  values should stack or reflow before truncating into ambiguity
\item
  blocks should be downstream of doctrine, patterns, and components
\item
  notification work remains parked while the design-language bridge is made explicit
\end{itemize}

\section{Implementation outline}\label{record-169-implementation-outline}

\begin{enumerate}
\def\labelenumi{\arabic{enumi}.}
\tightlist
\item
  Move DL-002 into the active cycle area and enrich it with playbacks and invariants.
\item
  Write cycle tests that demand:

  \begin{itemize}
  \tightlist
  \item
    canonical pattern language in the design-system docs
  \item
    an explicit blocks document
  \item
    updated design-system read order
  \item
    a spawned follow-on backlog item
  \end{itemize}
\item
  Update the design-system docs so the canonical patterns are written down where builders will actually read them.
\item
  Close the cycle with a retrospective and next-step backlog item.
\end{enumerate}

\section{Tests to write first}\label{record-169-tests-to-write-first}

Under \texttt{t\allowbreak{}e\allowbreak{}s\allowbreak{}t\allowbreak{}s\allowbreak{}/\allowbreak{}c\allowbreak{}y\allowbreak{}c\allowbreak{}l\allowbreak{}e\allowbreak{}s\allowbreak{}/\allowbreak{}D\allowbreak{}L\allowbreak{}-\allowbreak{}0\allowbreak{}0\allowbreak{}2\allowbreak{}/}:

\begin{itemize}
\tightlist
\item
  the active cycle doc includes the required workflow sections
\item
  the design-system docs explicitly define canonical selection, active-region, spacing, and overflow patterns
\item
  the design-system docs explicitly define what a block is and name the first block candidates
\item
  the Design Language legend points at DL-002 as the latest completed cycle and spawns the next backlog item
\end{itemize}

\section{Risks}\label{record-169-risks}

\begin{itemize}
\tightlist
\item
  documenting patterns too vaguely to guide implementation reviews
\item
  duplicating doctrine without actually bridging it into builder-facing docs
\item
  overcommitting to blocks before the component/pattern vocabulary is stable
\end{itemize}

\section{Out of scope}\label{record-169-out-of-scope}

\begin{itemize}
\tightlist
\item
  implementation of a reusable blocks runtime
\item
  service-level notification work
\item
  localization adapters or other legend work
\end{itemize}

\section{Retrospective}\label{record-169-retrospective}

\subsection{What landed}\label{record-169-what-landed}

\begin{itemize}
\tightlist
\item
  explicit canonical-pattern guidance in \texttt{d\allowbreak{}o\allowbreak{}c\allowbreak{}s\allowbreak{}/\allowbreak{}d\allowbreak{}e\allowbreak{}s\allowbreak{}i\allowbreak{}g\allowbreak{}n\allowbreak{}-\allowbreak{}s\allowbreak{}y\allowbreak{}s\allowbreak{}t\allowbreak{}e\allowbreak{}m\allowbreak{}/\allowbreak{}p\allowbreak{}a\allowbreak{}t\allowbreak{}t\allowbreak{}e\allowbreak{}r\allowbreak{}n\allowbreak{}s\allowbreak{}.\allowbreak{}m\allowbreak{}d} for:

  \begin{itemize}
  \tightlist
  \item
    selected rows
  \item
    active regions
  \item
    truthful visible controls
  \item
    spacing rhythm
  \item
    value-overflow behavior
  \end{itemize}
\item
  a new \texttt{d\allowbreak{}o\allowbreak{}c\allowbreak{}s\allowbreak{}/\allowbreak{}d\allowbreak{}e\allowbreak{}s\allowbreak{}i\allowbreak{}g\allowbreak{}n\allowbreak{}-\allowbreak{}s\allowbreak{}y\allowbreak{}s\allowbreak{}t\allowbreak{}e\allowbreak{}m\allowbreak{}/\allowbreak{}b\allowbreak{}l\allowbreak{}o\allowbreak{}c\allowbreak{}k\allowbreak{}s\allowbreak{}.\allowbreak{}m\allowbreak{}d} document that defines what a block is, where it sits in the design-system stack, and which first candidates are credible
\item
  updated design-system read order in \texttt{d\allowbreak{}o\allowbreak{}c\allowbreak{}s\allowbreak{}/\allowbreak{}d\allowbreak{}e\allowbreak{}s\allowbreak{}i\allowbreak{}g\allowbreak{}n\allowbreak{}-\allowbreak{}s\allowbreak{}y\allowbreak{}s\allowbreak{}t\allowbreak{}e\allowbreak{}m\allowbreak{}/\allowbreak{}R\allowbreak{}E\allowbreak{}A\allowbreak{}D\allowbreak{}M\allowbreak{}E\allowbreak{}.\allowbreak{}m\allowbreak{}d}
\end{itemize}

\subsection{Drift from ideal}\label{record-169-drift-from-ideal}

No material drift.

This cycle stayed documentation-first on purpose. It names the canonical patterns and block layer clearly, but it does not yet prove those rules in more shared implementation surfaces.

\subsection{Debt spawned}\label{record-169-debt-spawned}

Spawned:

\begin{itemize}
\tightlist
\item
  \href{https://github.com/flyingrobots/bijou/blob/cbd70612140e7a450d764b3c7be1c6af878b99bc/docs/method/graveyard/legacy-backlog/DL-003-prove-canonical-patterns-in-shared-surfaces.md}{DL-003 --- Prove Canonical Patterns in Shared Surfaces}
\end{itemize}

